
\setcounter{footnote}{20}

\medskip
\noindent\SGAInlineRefTarget{sga3:iii:remark:0-3}\textbf{Remark 0.3.}
Suppose that the \(X^+\)-scheme \(Y\) is flat over \(S\) (cf. SGA~1,
IV). We may then write
\[
 \operatorname{Hom}_{X^+}(Y,L_X)
 =
 \operatorname{Hom}_{\mathscr O_{Y_0}}
 \bigl(g_0^*(\Omega^1_{X_0/S_0}),
       \mathscr J\otimes_{\mathscr O_{S_0}}\mathscr O_{Y_0}\bigr).
\]

\medskip
\noindent\SGAInlineRefTarget{sga3:iii:remark:0-4}\textbf{Remark 0.4.}
Let \(\pi_0:X_0\to S_0\) be the structure morphism, and suppose that
there is an \(\mathscr O_{S_0}\)-module \(\omega^1_{X_0/S_0}\) such
that
\[
 \Omega^1_{X_0/S_0}\simeq
 \pi_0^*(\omega^1_{X_0/S_0}).
\]
(This occurs, in particular, when \(X_0\) is an \(S_0\)-group; cf.
\SGARefLink{sga3:ii:subsection:4-11}{II~4.11}.) Define a functor \(L'_X\) over \(S\) by
\begin{equation}
 \operatorname{Hom}_S(Y,L'_X)
 =
 \operatorname{Hom}_{\mathscr O_{Y_0}}
 \bigl(\omega^1_{X_0/S_0}\otimes_{\mathscr O_{S_0}}\mathscr O_{Y_0},
       \mathscr J\mathscr O_Y\bigr).
 \tag{0.4.1}\label{sga3:iii:remark:0-4:equation:0-4-1}
\end{equation}
Then
\[
 \operatorname{Hom}_{X^+}(Y,L_X)=\operatorname{Hom}_S(Y,L'_X)
\]
for every \(X^+\)-scheme \(Y\); in other words,
\[
 L_X=L'_X\mathbin{\times_S}X^+.
\]

Since
\(L_X\mathbin{\times_{X^+}}X=L'_X\mathbin{\times_S}X\), the action
of \(L_X\) on \(X\) induces an action of \(L'_X\) on \(X\), and
this action respects \(p_X:X\to X^+\).\footnote{Editor's note: what
follows has been expanded.} Indeed, if \(Y\) is an \(S\)-scheme,
\(h:Y\to X\) is an \(S\)-morphism, and \(m\in L'_X(Y)\), then
\(h\) and \(m\cdot h\) have the same restriction to \(Y_{\mathscr J}\);
that is,
\[
 p_X(m\cdot h)=p_X(h).
\]

\medskip
\noindent\SGAInlineRefTarget{sga3:iii:remark:0-5}\textbf{Remark 0.5.}
Keep the hypotheses and notation of \SGARefLink{sga3:iii:remark:0-4}{0.4}, and suppose in addition that
\(Y\) is an \(S\)-scheme flat over \(S\). Then
\[
 \operatorname{Hom}_{X^+}(Y,L_X)
 =\operatorname{Hom}_S(Y,L'_X)
 =\operatorname{Hom}_{S_0}(Y_0,L_{0X}),
\]
where the \(S_0\)-functor in abelian groups \(L_{0X}\) is defined by
the following identity (with \(T\) the variable \(S_0\)-scheme):
\begin{equation}
 \operatorname{Hom}_{S_0}(T,L_{0X})
 =
 \operatorname{Hom}_{\mathscr O_T}
 \bigl(\omega^1_{X_0/S_0}\otimes_{\mathscr O_{S_0}}\mathscr O_T,
       \mathscr J\otimes_{\mathscr O_{S_0}}\mathscr O_T\bigr).
 \tag{0.5.1}\label{sga3:iii:remark:0-5:equation:0-5-1}
\end{equation}
Thus, in the notation of II~1, we have shown that \(L'_X\) and
\(\prod_{S_0/S}L_{0X}\) have the same restriction to the full
subcategory of \((\mathbf{Sch})/S\) whose objects are the \(S\)-schemes
flat over \(S\).

\medskip
\noindent\SGAInlineRefTarget{sga3:iii:remark:0-6}\textbf{Remark 0.6.}
Keep the hypotheses and notation of \SGARefLink{sga3:iii:remark:0-5}{0.5},\footnote{Editor's note: the
following hypothesis, which was implicit in the original, has been
added.} and suppose in addition that \(\pi_0:X_0\to S_0\) has a
section \(\varepsilon_0\). Then
\[
 \omega^1_{X_0/S_0}\simeq
 \varepsilon_0^*(\Omega^1_{X_0/S_0}).
\]
First, independently of this last hypothesis, one has
\[
 \operatorname{Hom}_{S_0}(T,L_{0X})
 =
 \Gamma\!\left(
   T,
   \mathscr{H}om_{\mathscr O_T}
   \bigl(\omega^1_{X_0/S_0}\otimes_{\mathscr O_{S_0}}\mathscr O_T,
         \mathscr J\otimes_{\mathscr O_{S_0}}\mathscr O_T\bigr)
 \right).
\]

Suppose now that \(\omega^1_{X_0/S_0}\) admits a finite presentation
(cf. EGA~0\textsubscript{I}, 5.2.5), as is in particular the case when
\(X_0\) is locally of finite presentation over \(S_0\) (cf.
EGA~IV\textsubscript{4}, 16.4.22). If \(T\) is flat over \(S_0\), it
then follows from EGA~0\textsubscript{I}, 6.7.6 that
\[
 \mathscr{H}om_{\mathscr O_T}
 \bigl(\omega^1_{X_0/S_0}\otimes_{\mathscr O_{S_0}}\mathscr O_T,
       \mathscr J\otimes_{\mathscr O_{S_0}}\mathscr O_T\bigr)
 \simeq
 \mathscr{H}om_{\mathscr O_{S_0}}
 \bigl(\omega^1_{X_0/S_0},\mathscr J\bigr)
 \otimes_{\mathscr O_{S_0}}\mathscr O_T,
\]
whence
\[
 \operatorname{Hom}_{S_0}(T,L_{0X})
 =
 \Gamma\!\left(
   T,
   \mathscr{H}om_{\mathscr O_{S_0}}
   \bigl(\omega^1_{X_0/S_0},\mathscr J\bigr)
   \otimes_{\mathscr O_{S_0}}\mathscr O_T
 \right).
\]
Using the notation \(\mathbf W(\,\cdot\,)\) of \SGARefLink{sga3:i:definition:4-6-1}{I~4.6.1}, we have
therefore proved that, for every \(S_0\)-scheme \(T\) flat over
\(S_0\),
\[
 \operatorname{Hom}_{S_0}(T,L_{0X})
 =
 \operatorname{Hom}_{S_0}
 \left(T,
 \mathbf W\!\left(
   \mathscr{H}om_{\mathscr O_{S_0}}
   (\omega^1_{X_0/S_0},\mathscr J)
 \right)\right).
\]

In summary, if \(\omega^1_{X_0/S_0}\) admits a finite presentation
and one restricts to the category of \(S\)-schemes flat over \(S\),
then
\begin{equation}
 L'_X
 =
 \prod_{S_0/S}
 \mathbf W\!\left(
   \mathscr{H}om_{\mathscr O_{S_0}}
   (\omega^1_{X_0/S_0},\mathscr J)
 \right),
 \tag{0.6.1}\label{sga3:iii:remark:0-6:equation:0-6-1}
\end{equation}
and
\(\mathscr{H}om_{\mathscr O_{S_0}}
(\omega^1_{X_0/S_0},\mathscr J)\) is a quasi-coherent
\(\mathscr O_{S_0}\)-module, by EGA~I, 9.1.1.

Finally, suppose moreover that \(\omega^1_{X_0/S_0}\) is locally free
of finite rank, for example that \(X_0\) is smooth over \(S_0\) (in
which case it is automatically locally of finite presentation over
\(S_0\)). Then
\begin{equation}
 \mathscr{H}om_{\mathscr O_{S_0}}
   (\omega^1_{X_0/S_0},\mathscr J)
 \simeq
 \underline{\operatorname{Lie}}(X_0/S_0)
 \otimes_{\mathscr O_{S_0}}\mathscr J,
 \tag{0.6.2}\label{sga3:iii:remark:0-6:equation:0-6-2}
\end{equation}
where, by an abuse of language (since \(X_0\) need not be an
\(S_0\)-group),
\(\underline{\operatorname{Lie}}(X_0/S_0)\) denotes the dual of the
\(\mathscr O_{S_0}\)-module \(\omega^1_{X_0/S_0}\).\footnote{Editor's
note: this is justified by \SGARefLink{sga3:ii:proposition:3-3}{II~3.3} and \SGARefLink{sga3:iii:paragraph:4-11}{4.11}. Indeed, the functor
\(L^{\varepsilon_0}_{X_0/S_0}\), the ``tangent space to \(X_0\) over
\(S_0\) at \(\varepsilon_0\),'' is represented over \(S_0\) by the
vector bundle
\(\mathbf V(\varepsilon_0^*(\Omega^1_{X_0/S_0}))
=\mathbf V(\omega^1_{X_0/S_0})\). Its sheaf of sections is the dual
module
\(\mathscr{H}om_{\mathscr O_{S_0}}
(\omega^1_{X_0/S_0},\mathscr O_{S_0})\); this is
\(\underline{\operatorname{Lie}}(X_0/S_0)\) when \(X_0\) is an
\(S_0\)-group and \(\varepsilon_0\) is the identity section.}

\SGARefLink{sga3:iii:proposition:0-2}{Proposition~0.2} (and its proof) has two important
corollaries.\footnote{Editor's note: \SGARefLink{sga3:iii:corollary:0-7-bis}{Corollary 0.7 bis}, which was used
implicitly in the proof of \SGARefLink{sga3:iii:proposition:0-8}{0.8}, has been added to \SGARefLink{sga3:iii:corollary:0-7}{0.7}. We point out
that \SGARefLink{sga3:iii:proposition:0-8}{0.8}--\SGARefLink{sga3:iii:corollary:0-12}{0.12}, as well as \SGARefLink{sga3:iii:corollary:0-17}{0.17}, which play an important technical role
later in this expos\'e, are consequences of \SGARefLink{sga3:iii:corollary:0-7}{0.7} and \SGARefLink{sga3:iii:corollary:0-7-bis}{0.7 bis}.}

\medskip
\noindent\SGAInlineRefTarget{sga3:iii:corollary:0-7}\textbf{Corollary 0.7.}
\textit{Let \(X\) be an \(S\)-scheme.}
\begin{enumerate}
 \item[\SGAInlineRefTarget{sga3:iii:corollary:0-7:item:a}(a)] \textit{Every \(S\)-endomorphism of \(X\) inducing the
 identity on \(X_{\mathscr J}\) is an automorphism.}

 \item[\SGAInlineRefTarget{sga3:iii:corollary:0-7:item:b}(b)] \textit{There is an exact sequence of groups}
 \[
  0\longrightarrow
  \operatorname{Hom}_{\mathscr O_{X_0}}
    (\Omega^1_{X_0/S_0},\mathscr J\mathscr O_X)
  \xrightarrow{\ i\ }
  \operatorname{Aut}_S(X)
  \longrightarrow
  \operatorname{Aut}_{S_{\mathscr J}}(X_{\mathscr J}).
 \]

 \item[\SGAInlineRefTarget{sga3:iii:corollary:0-7:item:c}(c)] \textit{Let \(\operatorname{Aut}_S(X)\) act on the first
 group by transport of structure. Then, for every
 \(u\in\operatorname{Aut}_S(X)\) and
 \(m\in\operatorname{Hom}_{\mathscr O_{X_0}}
 (\Omega^1_{X_0/S_0},\mathscr J\mathscr O_X)\),}
 \[
  i(u\cdot m)=u\,i(m)\,u^{-1}.
 \]
\end{enumerate}

\noindent\textit{Proof.}
By \SGARefLink{sga3:iii:proposition:0-2:item:i}{0.2(i)}, \(\operatorname{Hom}_{X^+}(X,X)\) is a principal
homogeneous set under \(\operatorname{Hom}_{X^+}(X,L_X)\), since it is
certainly nonempty: it contains the marked point given by the identity
automorphism of \(X\).\footnote{Editor's note: what follows has been
expanded.} Consequently, the map \(m\mapsto m\cdot\operatorname{id}_X\)
induces a bijection
\[
 \operatorname{Hom}_{\mathscr O_{X_0}}
   (\Omega^1_{X_0/S_0},\mathscr J\mathscr O_X)
 =L_X(X)
 \xrightarrow{\sim}
 \operatorname{Hom}_{X^+}(X,X).
\]

Let \(m\in L_X(X)\), and let
\(f=m'\cdot\operatorname{id}_X\) be an element of
\(\operatorname{Hom}_{X^+}(X,X)\). Applying \SGARefLink{sga3:iii:proposition:0-2:item:ii}{0.2(ii)} to \(f\) gives
\[
 f\circ(m\cdot\operatorname{id}_X)
 =L_f(X)(m)\cdot f
 =L_f(X)(m)\cdot(m'\cdot\operatorname{id}_X).
\]
On the other hand, because \(f\) is an \(X^+\)-endomorphism of \(X\),
one has \(f_{\mathscr J}=\operatorname{id}_{X_{\mathscr J}}\), hence
\(f_0=\operatorname{id}_{X_0}\). Since \(L_f\) depends only on
\(f_0\) (cf. the editor's note (17) in \SGARefLink{sga3:iii:corollary:0-1-10}{0.1.10}), it follows that
\(L_f(X)(m)=m\). Thus the preceding equality becomes
\[
 (m'\cdot\operatorname{id}_X)\circ
 (m\cdot\operatorname{id}_X)
 =m\cdot(m'\cdot\operatorname{id}_X)
 =(m+m')\cdot\operatorname{id}_X.
\]
This shows that the bijection
\(m\mapsto m\cdot\operatorname{id}_X\) carries the group law of
\(\operatorname{Hom}_{X^+}(X,L_X)\) to composition of
\(X^+\)-endomorphisms of \(X\).

It follows first that every element of
\(\operatorname{Hom}_{X^+}(X,X)\) is invertible, proving part~\SGARefLink{sga3:iii:corollary:0-7:item:a}{(a)}, and
then that there is an exact sequence
\[
 0\longrightarrow \operatorname{Hom}_{X^+}(X,L_X)
 \xrightarrow{\ i\ }\operatorname{Aut}_S(X)
 \longrightarrow\operatorname{Aut}_{S_{\mathscr J}}(X_{\mathscr J}),
\]
which proves part~\SGARefLink{sga3:iii:corollary:0-7:item:b}{(b)}.

The morphism \(i\) just defined is functorial in \(X\) with respect to
isomorphisms, because it is defined structurally from the action of
\(L_X\) on \(X\) over \(X^+\), an action that is itself functorial in
\(X\) by \SGARefLink{sga3:iii:proposition:0-2:item:ii}{Proposition~0.2(ii)}.\footnote{Editor's note: what follows has
been expanded.} Hence every automorphism \(u\) of \(X\) over \(S\)
induces, by transport of structure, isomorphisms
\[
 h:\operatorname{Hom}_{X^+}(X,L_X)
   \xrightarrow{\sim}\operatorname{Hom}_{X^+}(X,L_X)
 \quad\text{and}\quad
 f:\operatorname{Aut}_S(X)
   \xrightarrow{\sim}\operatorname{Aut}_S(X)
\]
such that the following diagram commutes.

\SGAThreeNativeDiagram[0.70\textwidth]
  {native_diagrams/exp3/diagram_022_exp3_p016_cor07_functoriality_square.tex}
  {Functoriality of \(i\) under transport of structure.}
  {fig:exp3-cor07-functoriality}
  {Polo--Gille Exp3-14oct24.pdf, component-local p. 16; combined-reader p. 126; printed p. 116}

Thus \(f\circ i=i\circ h\). On the other hand, \(f\) is described by
the following commutative diagram.

\SGAThreeNativeDiagram[0.40\textwidth]
  {native_diagrams/exp3/diagram_023_exp3_p016_cor07_conjugation_square.tex}
  {The automorphism \(f\) acts by conjugation with \(u\).}
  {fig:exp3-cor07-conjugation}
  {Polo--Gille Exp3-14oct24.pdf, component-local p. 16; combined-reader p. 126; printed p. 116}

In other words,
\[
 f(a)=u\circ a\circ u^{-1}
 \qquad\text{for every }a\in\operatorname{Aut}_S(X).
\]
Writing \(i(h(m))=f(i(m))\) now gives the asserted formula.

